\subsection{A fixed incoming arrow and the extreme eight-element corner}

\begin{lemma}[A fixed incoming arrow]
\label{C:fixed-incoming-cups}
Let $C$ be a finite category whose objects consist of a full subcategory
$D$ and one additional object $v$, with no arrows from $v$ to $D$.
Write $M=\operatorname{End}_C(v)$.  Suppose
$H^2(BD;\mathbb Q)=0$ and there is an arrow $f:d\to v$, for some
$d\in D$, such that $mf=f$ for every $m\in M$.
Then
\[
 H^2(BC;\mathbb Q)\otimes H^2(BC;\mathbb Q)
       \longrightarrow H^4(BC;\mathbb Q)
\]
is the zero map.  The opposite assertion holds for a source object
and a fixed outgoing arrow.
\end{lemma}
\begin{proof}
The correspondence $d\mapsto\operatorname{Hom}_C(d,v)$ has commuting
right $D$- and left $M$-actions.  Its two-sided bar space $E$ gives the
natural homotopy-pushout description
\[
 BC\simeq BD\mathop{\cup}\limits_E^h BM.
\]
This description follows directly by subdividing each mixed nerve simplex
at its unique crossing from $D$ to $v$; the pieces are the double bar,
with its two endpoint projections.

The chosen arrow gives an actual section $BM\to E$ of the projection
$E\to BM$.  In the double bar, keep the $D$-coordinate at the object $d$
and the correspondence-coordinate at $f$, and retain all the $M$-arrows.
The equalities $mf=f$ verify every endpoint face; the other faces and
degeneracies are the ordinary bar identities.  Thus this is a simplicial
section, rather than a merely linear splitting on homology.  Consequently
$H^2(BM;\mathbb Q)\to H^2(E;\mathbb Q)$ is injective.

Let $X$ be the double-mapping-cylinder model of the displayed pushout,
and let $A=BD\sqcup BM\subset X$ be its two endpoint spaces.
The Mayer--Vietoris sequence and $H^2(BD;\mathbb Q)=0$ show that
$H^2(X;\mathbb Q)\to H^2(A;\mathbb Q)$ is zero.  Every degree-two
class therefore lifts to $H^2(X,A;\mathbb Q)$.  The quotient is
\[
 X/A\cong (E\times[0,1])/(E\times\{0,1\})
       \cong \Sigma(E_+),
\]
where the entire endpoint subspace is collapsed to one basepoint.
Positive-degree products in its reduced cohomology vanish because it is
a suspension.  Relative products followed by the map to absolute
cohomology now prove the claim.
\end{proof}

\begin{corollary}[The extreme eight-element corner at fifteen arrows]
\label{C:fifteen-extreme-eight}
None of the following linear three-object cardinality profiles, nor
any of their opposites, can have type $\mathbb{CP}^2$:
\[
 (m,n,p,x,y,z)=(1,1,8,2,1,2),\quad
 (1,1,8,2,2,1),\quad (2,1,8,1,1,2).
\]
Here $(m,n,p)$ are the endomorphism orders at $0,1,2$ and
$(x,y,z)$ are the cardinalities of the hom-sets from $0$ to $1$,
from $1$ to $2$, and from $0$ to $2$, respectively.
No classification of the eight-element corner is required.
\end{corollary}
\begin{proof}
Take $D$ to be the full subcategory on $0,1$ and the sink to be $2$.
For the first two profiles, $BD$ is the circle formed by two parallel
arrows, so $H^2(BD;\mathbb Q)=0$.  In the first profile the singleton
hom-set from $1$ to $2$ supplies a fixed incoming arrow.  In the second,
the singleton hom-set from $0$ to $2$ supplies it.  For the third profile,
object $1$ is terminal in $D$, because its endomorphism monoid and the
hom-set from $0$ to $1$ are singletons.  Hence $BD$ is contractible, and
the singleton hom-set from $1$ to $2$ again supplies the fixed arrow.
Lemma~\ref{C:fixed-incoming-cups} applies in all three cases, whereas
the rational degree-two generator of $\mathbb{CP}^2$ has nonzero square.
\end{proof}
